\section{Adaptations in \texttt{fss\_SAT.py}}
\label{sec:adaptations}

The changes fall into three groups: performance rewrites that leave the result
of \citet{faggian2015} unchanged (Sections~\ref{sec:filter}--\ref{sec:score}),
feature extensions (Sections~\ref{sec:modes}--\ref{sec:parallel}), and
algorithmic additions that go beyond the original method
(Sections~\ref{sec:nan}--\ref{sec:cwfss}). Table~\ref{tab:overview} gives an
overview.

\begin{table}[ht!]
\centering
\renewcommand{\arraystretch}{1.25}
\begin{tabularx}{\textwidth}{l l X}
\toprule
\textbf{Where} & \textbf{Change} & \textbf{Effect} \\
\midrule
\texttt{integral\_filter}   & \texttt{np.pad} + slicing   & same result, faster, no fancy indexing \\
\texttt{\_fss\_score}        & BLAS dot products           & same result, no temporary arrays \\
\texttt{\_build\_binary\_sat} & threshold modes            & \texttt{over}/\texttt{under}/\texttt{between}/\texttt{tolerance} \\
\texttt{fss\_threshold}      & percentile thresholds       & threshold given as a percentile of the field \\
\texttt{fss\_threshold}      & over-estimation \texttt{ovest} & frequency-bias diagnostic returned with FSS \\
\texttt{fss\_threshold\_eps} & ensemble FSS                & 3D forecast stack, exceedance probability \\
\texttt{fss\_cumsum\_parallel} & \texttt{joblib}           & threshold sweep optionally parallel (\texttt{n\_jobs}) \\
\texttt{\_fss\_score\_masked} & per-window valid weighting & transparent missing-data (NaN) handling \\
\texttt{CWFSS}               & R2 sampling, bootstrap, NaN & continuous/weighted FSS \\
\bottomrule
\end{tabularx}
\caption{Overview of the adaptations made on top of the \citet{faggian2015}
reference implementation.}
\label{tab:overview}
\end{table}

\subsection{Vectorised neighbourhood filter}
\label{sec:filter}

The reference \codeword{integral_filter} (Section~\ref{sec:baseline}) builds two
full coordinate grids with \codeword{np.mgrid}, clips them, and gathers the four
corner tables with fancy indexing. The rewrite achieves the identical clamped
box-sum of Eq.~\eqref{eq:boxsum} with edge-replicated padding and plain
contiguous slicing, which avoids the large index arrays and the gather:

{\setstretch{1}\begin{lstlisting}[style=python]
w = n // 2
if w < 1:
    return field
rows, cols = field.shape[-2], field.shape[-1]
D = 2 * w
p = np.pad(field, w, mode='edge')          # 2D case
return (p[D:D+rows, D:D+cols] + p[:rows, :cols]
      - p[:rows, D:D+cols] - p[D:D+rows, :cols])
\end{lstlisting}}

Padding the integral table with \codeword{mode='edge'} replicates the boundary
rows and columns, so a corner lookup that would have been clamped to the last
valid index instead reads the replicated edge value --- exactly what the
original \codeword{np.clip} produced. The boundary behaviour is therefore
\emph{identical} to the reference implementation; only the data access pattern
changes. The 3D (ensemble) case pads only the two trailing spatial axes.

\subsection{BLAS dot-product score}
\label{sec:score}

The reference \codeword{fss} scales both windowed fields to fractions and then
forms two temporary arrays, \codeword{(fhat-ohat)**2} and
\codeword{fhat**2+ohat**2}, before averaging. The rewrite (\codeword{_fss_score})
expresses the same quantities through three inner products of the
\emph{unscaled} box sums $S_f$, $S_o$, which BLAS evaluates without allocating
temporaries:
%
\begin{equation}
  \begin{aligned}
  \mathit{ff} &= \textstyle\sum_i S_f^2, &\quad
  \mathit{oo} &= \textstyle\sum_i S_o^2, &\quad
  \mathit{fo} &= \textstyle\sum_i S_f S_o .
  \end{aligned}
\end{equation}
%
Using the identity $\sum_i (S_f-S_o)^2 = \mathit{ff} + \mathit{oo} -
2\,\mathit{fo}$ and writing $A=(2w+1)^2$ for the window area, the numerator and
denominator of Eq.~\eqref{eq:fss} become
%
\begin{equation}
  \mathrm{num} = \frac{1}{A^2}\,\frac{\mathit{ff}+\mathit{oo}-2\,\mathit{fo}}{N},
  \qquad
  \mathrm{denom} = \frac{1}{A^2}\,\frac{\mathit{ff}+\mathit{oo}}{N},
\end{equation}
%
where $A^{-2}$ is the constant \codeword{inv_area_sq} factor (the fractions
$p=S/A$ enter squared). This is algebraically identical to the reference
result, but the score reduces to three \codeword{np.dot} calls plus scalar
arithmetic:

{\setstretch{1}\begin{lstlisting}[style=python]
ff = np.dot(fflat, fflat); oo = np.dot(oflat, oflat); fo = np.dot(fflat, oflat)
num   = inv_area_sq * (ff + oo - 2.0 * fo) / n
denom = inv_area_sq * (ff + oo) / n
if denom == 0.0:
    return 0.0, 0.0, np.nan
return num, denom, 1.0 - num / denom
\end{lstlisting}}

Note that \codeword{np.dot} is not NaN-aware, unlike the reference
\codeword{np.nanmean}. This fast path is therefore only taken when the fields
contain no missing data; the missing-data case is handled separately
(Section~\ref{sec:nan}).

\subsection{Threshold modes and percentile thresholds}
\label{sec:modes}

The reference implementation binarises with a single ``over'' comparison
($\text{field} > t$). \texttt{fss\_SAT.py} factors the binarisation into
\codeword{_build_binary_sat} and supports four \codeword{threshold_mode} values:
%
\begin{description}
  \item[\texttt{over}] $\text{field} > t$ (default, the classic FSS);
  \item[\texttt{under}] $\text{field} \le t$;
  \item[\texttt{between}] $t_1 < \text{field} \le t_2$, for verification of a
        value band (in the sweep, an extra lower edge $-1$ is inserted so that
        the lowest band still captures the zero values);
  \item[\texttt{tolerance}] $(1-\tau)\,t < \text{field} \le (1+\tau)\,t$,
        a relative tolerance band of half-width $\tau$ (\codeword{tolerance})
        around $t$.
\end{description}
%
Independently, \codeword{percentiles=True} interprets the threshold as a
percentile of the field rather than an absolute value; the forecast and
observation thresholds are then derived separately with \codeword{np.percentile}
(\codeword{np.nanpercentile} on the missing-data path).

\subsection{Over-estimation diagnostic}
\label{sec:ovest}

Alongside the numerator, denominator and score, \codeword{fss_threshold} now
returns an over-estimation (frequency-bias) value
%
\begin{equation}
  \mathrm{ovest} = \frac{\#\{\text{fcst} > t\} - \#\{\text{obs} > t\}}
                        {N_\text{points}},
\end{equation}
%
the difference between the forecast and observed event counts normalised by the
number of points. It is a domain-wide bias indicator and is broadcast to the
shape of the window axis so it travels through the same return structure as the
score. On the missing-data path it is computed over valid points only.

\subsection{Ensemble FSS (eFSS)}
\label{sec:eps}

\codeword{fss_threshold_eps} implements the ensemble extension of the FSS in the
spirit of \citet{duc2013}. The forecast is a 3D stack (members $\times$ rows
$\times$ cols). Instead of a binary forecast field, the per-point
\emph{exceedance probability} across members is used,
%
\begin{equation}
  p_f(i,j) = \frac{1}{M}\sum_{m=1}^{M}\mathbb{1}\!\left[\,\text{fcst}_m(i,j)
             \text{ exceeds } t\,\right],
\end{equation}
%
($M$ = number of members), and this probability field is fed into the same SAT
machinery as the deterministic case. The observation remains a single 2D field.
An assertion enforces the 3D forecast shape.

\subsection{Threshold sweep and parallelisation}
\label{sec:parallel}

\codeword{fss_cumsum_parallel} sweeps a list of thresholds, calling
\codeword{fss_threshold} (or \codeword{fss_threshold_eps} when
\codeword{eps=True}) once per threshold and reusing the cached integral tables
across all windows. The sweep can run serially (\codeword{n_jobs=1}) or in
parallel over thresholds through \codeword{joblib} (\codeword{n_jobs>1}):

{\setstretch{1}\begin{lstlisting}[style=python]
if n_jobs == 1:
    ret = [use_fss_threshold_func(fcst, obs, t1, t2, windows, ...) for t1, t2 in calls]
else:
    from joblib import Parallel, delayed
    ret = Parallel(n_jobs=n_jobs)(
        delayed(use_fss_threshold_func)(fcst, obs, t1, t2, windows, ...) for t1, t2 in calls)
\end{lstlisting}}

\codeword{fss_cumsum_frame} wraps the sweep and returns pandas
\codeword{DataFrame}s of numerator, denominator, score and over-estimation
(indexed by threshold, columns by window), or just the raw score array when
\codeword{raw=True}. Note that, because BLAS itself is multi-threaded,
oversubscription can occur when \codeword{n_jobs>1}; the benchmark
\codeword{bench_njobs.py} measures this with and without pinning the BLAS thread
count.

\subsection{Missing-data handling via per-window valid-point weighting}
\label{sec:nan}

This is the most substantial addition. The reference method has no concept of
missing data; a single NaN would propagate through the cumulative sums and
poison the whole table. \texttt{fss\_SAT.py} adds a mask-aware path that treats
a grid point as \emph{missing} when either the forecast \emph{or} the
observation is NaN there (\codeword{_validity_mask}). The clean path is taken
whenever the mask is all-valid, so fields without NaNs are unaffected and pay no
extra cost.

When NaNs are present, the binary fields are zeroed at missing points before the
SAT is built (\codeword{_build_binary_sat_masked}), so missing points never
contribute to a window sum. Each window is then weighted by its number of
\emph{valid} points,
%
\begin{equation}
  C = A - (\text{missing points in the window}),
  \label{eq:validcount}
\end{equation}
%
where $A=(2w+1)^2$ is the full window area and the missing-point count itself is
obtained from a SAT of the invalid mask (\codeword{_invalid_sat}) sampled with
the same \codeword{integral_filter}. With $S_f$, $S_o$ the box sums of the
mask-zeroed binary fields, the windowed fractions are $p_f=S_f/C$, $p_o=S_o/C$,
and weighting each window centre by $C$ collapses the weighted means to
%
\begin{equation}
  \mathrm{num} = \frac{\sum_w (S_f-S_o)^2 / C}{\sum_w C},
  \qquad
  \mathrm{denom} = \frac{\sum_w (S_f^2+S_o^2) / C}{\sum_w C}
  \label{eq:maskedscore}
\end{equation}
%
(\codeword{_fss_score_masked}). Windows with no valid points ($C=0$) drop out
naturally. The construction is deliberately consistent with the clean path: with
no missing data $C\equiv A$ everywhere, and Eq.~\eqref{eq:maskedscore} reduces
\emph{bit-for-bit} to the standard score of Section~\ref{sec:score}. Percentile
thresholds on this path use \codeword{np.nanpercentile}.

The same weighting is wired into the ensemble path (\codeword{fss_threshold_eps},
Section~\ref{sec:eps}): there a grid point counts when the observation is valid
\emph{and} at least one member is valid, and the exceedance probability is
averaged over the valid members only (\codeword{_eps_prob_masked}).

\subsection{Continuous/weighted FSS (\texttt{CWFSS})}
\label{sec:cwfss}

\codeword{CWFSS} aggregates the FSS over a \emph{range} of thresholds and
windows into a single scalar, instead of reporting the full
threshold\,$\times$\,window table. It draws \codeword{nsamples} (threshold,
window) pairs from the low-discrepancy R2 sequence (\codeword{R2}), which gives
even coverage of the 2D parameter space without a fixed grid, and computes the
FSS for each pair. Each sample is weighted by a threshold factor and a window
factor,
%
\begin{equation}
  \text{t\_factor} = \frac{t_\text{max}+t}{t_\text{max}},
  \qquad
  \text{w\_factor} = \frac{2\,w_\text{max}}{w_\text{max}+w},
\end{equation}
%
emphasising higher thresholds and smaller windows, and the weighted scores are
normalised by their achievable maximum to yield \codeword{self.cwfss}. A
\codeword{bootstrap} method resamples the stored per-sample scores to obtain a
confidence distribution.

\codeword{CWFSS} carries the same extensions as the batch functions. It honours
\codeword{threshold_mode}/\codeword{tolerance} through a shared
\codeword{_binarise} helper, and it is NaN-aware: threshold limits use
\codeword{np.nanmax}/\codeword{np.nanpercentile}, and when the fields contain
NaNs it uses the masked score of Eq.~\eqref{eq:maskedscore}. Samples whose FSS
is undefined (NaN, e.g.\ a window with no valid exceedance) are dropped from the
\codeword{np.nanmean}, and the theoretical maximum is masked identically so the
achievable ceiling stays at $1$ and is not lowered merely by the presence of
missing data.
